\documentclass[a4paper,11pt]{article}

\usepackage[margin=2.2cm]{geometry}
\usepackage[T1]{fontenc}
\usepackage[utf8]{inputenc}
\usepackage{xcolor}
\usepackage{enumitem}
\usepackage{paracol}

\definecolor{accent}{HTML}{7A2E2E}

\pagestyle{empty}

\newcommand{\recipetitle}[2]{%
  \begin{center}
    {\color{accent}\rule{0.4\linewidth}{1pt}}\\[4pt]
    {\Large\bfseries #1}\\[2pt]
    {\small\itshape #2}\\[2pt]
    {\color{accent}\rule{0.4\linewidth}{1pt}}
  \end{center}
  \vspace{0.6em}
}

\begin{document}

\recipetitle{Braised Lamb Shanks with Rosemary}{Serves 4 \quad|\quad Prep 20 min \quad|\quad Cook 2h 30min}

\begin{paracol}{2}
\setlength{\columnsep}{1.8em}

\begin{leftcolumn}
  {\color{accent}\bfseries Ingredients}
  \vspace{0.4em}

  \begin{itshape}
  \begin{itemize}[leftmargin=1.1em, itemsep=3pt, label={\color{accent}\textbullet}]
    \item 4 lamb shanks
    \item 2 tbsp olive oil
    \item 1 onion, diced
    \item 2 carrots, diced
    \item 2 celery stalks, diced
    \item 4 garlic cloves, crushed
    \item 250\,ml red wine
    \item 400\,g tinned chopped tomatoes
    \item 500\,ml lamb or beef stock
    \item 3 sprigs fresh rosemary
    \item 2 bay leaves
    \item Salt and black pepper
  \end{itemize}
  \end{itshape}

  \vspace{1em}
  {\color{accent}\bfseries Wine Pairing}
  \vspace{0.4em}

  \textit{A medium-bodied red such as a Chianti or Côtes du Rhône matches
  the dish's rosemary and tomato notes without overpowering the lamb.}
\end{leftcolumn}

\switchcolumn

  {\color{accent}\bfseries Method}
  \vspace{0.4em}

  \begin{enumerate}[leftmargin=1.6em, itemsep=6pt,
                     label={\color{accent}\bfseries\arabic*.}]
    \item Preheat the oven to 160\textdegree C (140\textdegree C fan).
      Season the lamb shanks generously with salt and pepper.
    \item Heat the olive oil in a large, oven-safe casserole dish over
      medium-high heat. Brown the shanks on all sides, about 3 minutes per
      side, then remove and set aside.
    \item Add the onion, carrots and celery to the same pot and cook for 5
      minutes until softened. Add the garlic and cook for 1 minute more.
    \item Pour in the red wine, scraping up any browned bits from the
      bottom of the pot. Simmer for 2--3 minutes to reduce slightly.
    \item Return the shanks to the pot. Add the chopped tomatoes, stock,
      rosemary and bay leaves. Bring to a gentle simmer.
    \item Cover and transfer to the oven. Braise for 2--2.5 hours, turning
      the shanks halfway through, until the meat is falling off the bone.
    \item Remove the shanks and keep warm. If the sauce is thin, reduce it
      on the stovetop over high heat for 5--10 minutes.
    \item Discard the bay leaves and rosemary stems. Return the shanks to
      the sauce and serve with mashed potatoes or polenta.
  \end{enumerate}

  \vspace{1em}
  {\color{accent}\bfseries Chef's Note}
  \vspace{0.4em}

  This dish improves overnight --- if time allows, cool and refrigerate
  after step 6, then reheat gently and finish with steps 7--8 the next day.

\end{paracol}

\end{document}
